
\documentclass[12pt]{report}

\usepackage{setspace}
\onehalfspacing %doublespace %

\usepackage{amssymb,amsmath,amsthm}

\setlength{\intextsep}{1ex}

%\usepackage[dvips]{graphicx}
\usepackage{graphicx}
\include{epstopdf}
\renewcommand{\qedsymbol}{\rule{0.4em}{0.4em}}
\usepackage{tabularx}

\usepackage[a4paper,bindingoffset=25mm,%
            left=0mm,right=15mm,top=15mm,bottom=15mm,%
            footskip=5mm]{geometry}

\usepackage{lipsum}  
 

\usepackage{float}
\setlength{\intextsep}{1ex}
\usepackage[shortlabels]{enumitem}
\usepackage{multirow}
\usepackage[font={small}]{caption}
\usepackage{nameref}

\usepackage[justification=centering]{caption}
\usepackage{caption,subcaption}
\usepackage[toc,page]{appendix}



\setlength{\parskip}{10pt}
\setlength{\parindent}{0pt}

\renewcommand{\bibname}{References}


\setlength\parindent{0pt}
\setlength{\unitlength}{1cm} 


\begin{document}

\begin{titlepage}
\centering

\begin{figure}[H]
\centering
\includegraphics[width=0.4\textwidth]{MAK_Logo2_png.png}%
\end{figure}


\textsc{}

\vspace{\stretch{0.1}}

{\huge \textsc{Design of a machine learning based system for pharmaceutical purchases}  \\}
\rule{3in}{0.4pt}

\vspace{\stretch{0.4}}

{\Large\textbf{Abraham Jerry Kakooza}	 \\}
Reg No.: 16/U/327\\ \vspace{2cm}
{\large \textbf{Department of Electrical and Computer Engineering} \\
School of Engineering \\
College of Engineering, Design, Art and Technology}



\vspace{\stretch{0.3}}

{
Supervisor: \textbf{Dr Andrew Katumba}    \\
Co-supervisor: \textbf{Prof. Eng. Dr. Peter Lating, iLabs Principal Investigator} }

\vspace*{\stretch{0.2}}

% Add date in desired format
April 27, 2020


\vspace*{\stretch{0.2}}
% Add the statement of what the report represents
{
A Report submitted in partial fulfillment of the requirements for the \\ \textbf{ Degree of Bachelor of Science in Telecommunications Engineering} \\ at Makerere University. }


\end{titlepage}



%\maketitle
				
%\beforeabstract %\input{abstract} 



\section*{Declaration}

Add your declaration here.



\newpage


\section*{Dedication}

\vspace{2cm}
\begin{center}
I dedicate this to my parents who have been there for me in this entire study period both emotionally and financially .
\end{center}



\newpage

\section*{Acknowledgments}

Add acknowledgments. 


\lipsum[2-4]



\newpage



\section*{Abstract}


Using a univariate time series analysis, we were able to validate different methods and approaches related to time series data preparation, analysis and forecast with an aim to facilitate recommending sales and marketing strategies based on trend/seasonality effects and forecasting sales of top five selling pharmaceutical products with diverse characteristics, such as stationarity, seasonality, amount of residuals and sales data variance. These drugs are Amoxicillin caps, Ampicloxa, Ceftriaxone Inj, Ciprofloxacin, Cotrimoxazole. 
......

\newpage

\tableofcontents

\newpage

\listoffigures

\newpage

\listoftables

\newpage



\section*{List of Abbreviations}

\begin{tabular}{p{3cm} l}
	ARIMA	& Auto Regression Integrated Moving Average\\
	LSTM	& Long Short Term Memory\\	
	

\end{tabular}


\newpage

\chapter{Introduction}

\section{Background}


People act on the basis of forecasting models whether they are on paper or in their heads. You are better off quantifying these estimations so you can discuss them rationally as opposed to making them based on intuition. For pharmaceutical distribution companies, it is essential to get good estimates of drugs due to the short shelf life of many medicines, the need to control stock levels and to avoid excessive inventory costs while guaranteeing customer demand satisfaction.This decreases the possibility of loss of customers due to stock outages.\\

Before digital technology dominated the world, the forecasting process was done manually by experienced individuals in the domain. This intuition required a lot of experience and was prone to error. Such a multifaceted problem started to realize the need for automating the pharmacy sales forecasting process. Research has been carried out with statistical, machine learning, deep learning and ensemble techniques to achieve more accurate sales forecasts. \cite{Perera2019}\\

\section{Aims and Objectives}

The main aim of the research project was to precisely predict sales of drugs and medical supplies.

The specific objectives  were:

\begin{itemize}[topsep=0pt]

\item To collect datasets of previous sales and purchases from Soteria Pharmacy.

\item To train and validate datasets with ARIMA and LSTM models.

\item To optimize the best model algorithm for accurate performance.

\end{itemize}













\section{Contributions of the Project}

The following major contributions have been accomplished during the course of this research:

\begin{enumerate}[topsep=0pt]

\item Developed a robust neural network model that could accurately predict future sales for a period of 50 days using a 2 year period worth of sales and purchases information.

\item Performed a data cleaning task, analysis and training with the various Machine learning models and came up with a graphical output.

\item Proposed a multivariate input approach with other characteristics such as precipitational weather information and promotional sales.

\end{enumerate}






\section{Organization of the Report}

This report has 3 chapters. 

\begin{itemize}
	
	\item   Introduction
	\item   Analysis and Findings
	\item 	Conclusion
\end{itemize}

% Describe what each chapter discusses.















\chapter{Heading of Chapter 2}
% Usually Literature Review

\section{Introduction}

% Start with an introduction


\section{....}

% Add sections and (sub)subsections as appropriate


\lipsum[2-4]






\section{Summary}
% End your chapter with a short summary.































\chapter{Add another Chapter}

\section{Introduction}

% Short Introduction

\section{Mathematical Preliminaries}

% Add sections and (sub)subsections as appropriate


\lipsum[2-4]


% Add other chapters and their (sub)(sub)sections as appropriate







\chapter{Numerical Results}

% Usually, start with a table of parameters similar to the one below, if this applies to your project.


\begin{table}[H]
\caption{Simulation Parameters} % title of Table
\centering % used for centering table
    \begin{tabular}{| l | l |}
	 \hline\hline
    Parameters & Values \\ \hline
    Network size &  $A = 5\,\text{km}\times 5\,\text{km}$ \\ \hline
    Total BS bandwidth &  $\mathcal{B}=20\,\text{MHz}$ \\ \hline
		BS and user densities &  $\lambda_b = 1.6\times 10^{-5}\,\,\text{m}^2$, $\lambda_u = 2\lambda_b$  \\ \hline
    Transmit power &  $P_t = 21\, \text{dBm}$ \\ \hline
    Pathloss parameters & $L = -33\,\text{dB}$, $\alpha=4$ \\ \hline
    Power parameters & $\mathcal{N}_b = 2$, $P_0=6.8\,\text{W}$, $\Delta=4$, $P_{sl}=4.3$ \\\hline	
    Additive noise parameters & $F = 10$, $T_a = 300\,\text{K}$\\\hline		
\end{tabular}
%\label{tab:parameters-sleep-mode} 
\end{table}



\lipsum[2-4]



% End with a conclusion chapter
\chapter{Conclusions and Future Works}
%\label{chap:conclusion}

\section{Conclusions}

%Conclude here


\lipsum[2-4]



\section{Ideas for Future Work}

% Add any ideas for further work here


% You can also add a section on challenges and recommendations if you wish












%\bibliographystyle{alpha}




\bibliographystyle{plain}
\bibliography{doc}




%% the bibliography style determines the format  in which both citations and references are printed,
%% other possible values are plain and abbrv
%%
%% If you want more control of the format of your citations you might want to take a look at
%% natbib.sty, which should be part of any standard LaTeX installation
%%
%% University regulations simply require that your citation style be consistent, so see what style
%% your supervisor recommends.

% Appendices start here



\appendix

\chapter{Appendix}

\lipsum[2-3]

\end{document}
